This is a generic appendix. Obcaecati cupiditate non provident, similique sunt in culpa, qui officia deserunt mollitia animi, id est laborum\index{laborum} et dolorum fuga. Duis aute irure dolor in reprehenderit in voluptate\index{voluptate} velit esse cillum dolore eu fugiat nulla pariatur. Excepteur sint occaecat cupidatat non proident, sunt in culpa qui officia deserunt mollit anim id est laborum. 

\section{Generic appendix section} 

This is a generic subsection in the appendix. Lorem ipsum dolor sit amet, consectetur adipisicing elit, sed do eiusmod tempor incididunt ut labore et dolore magna aliqua. Ut enim ad minim veniam, quis nostrud exercitation ullamco laboris nisi ut aliquip ex ea commodo consequat\index{consequat}. 

In efficitur mi sed eros tincidunt tempus. Cras ut neque vehicula, convallis odio in, lacinia orci. Vivamus et nunc vestibulum, efficitur eros vitae, bibendum justo. Aenean tincidunt ligula ut augue pulvinar, ut placerat risus accumsan. Integer pellentesque odio ac commodo dictum. Ut ut justo a mi consequat vestibulum. 

Lorem ipsum dolor sit amet, consectetur adipiscing elit. Fusce gravida aliquam velit, sed fermentum\index{fermentum} odio fermentum eu. Donec eget fermentum ligula. Nulla facilisi. Integer vitae metus eget risus vulputate consectetur. Suspendisse potenti. Quisque dictum purus non justo aliquet, sed tristique odio dapibus. Sed tincidunt justo vitae leo mattis, a ullamcorper metus vestibulum. Sed ac sapien metus. Sed non mauris nec nisl vehicula iaculis eget eget enim. Suspendisse potenti. Sed ac purus volutpat, venenatis sem eget, fermentum lorem. Suspendisse eu convallis tortor. Suspendisse potenti. Nunc vel dictum libero. Cras eleifend, orci non faucibus gravida, sapien quam pharetra ligula, nec tincidunt elit ex at sapien. 

\section{Assumptions and theorems and proofs}

This section displays an assumption and results to illustrate how they are typeset and numbered in the appendix. Note in particular that numbering of all the results is specific to each appendix (just as it is specific to each chapter in the main matter). First is an assumption:
\begin{assumption} 
Similique sunt in culpa, qui officia deserunt mollitia animi, id est laborum et dolorum fuga:
\begin{equation*}
\mathbb{E}(\Omega_{m,n}) = \mathbb{P}(\omega\cdot \mu - \xi) - \sum_{i=0}^{m}\sum_{j=1}^{n} \sigma(i,j) + 123^{56}.
\end{equation*}
Curabitur volutpat ultrices nunc id efficitur. Donec posuere, dui vel gravida elementum, ligula nisi lobortis elit, eu aliquet quam ligula id libero. Sed ac efficitur tur.
\label{a:appe}\end{assumption}

Lorem ipsum dolor sit amet, consectetur adipisicing elit, sed do eiusmod
tempor incididunt ut labore et dolore magna aliqua. Ut enim ad minim veniam,
quis nostrud exercitation ullamco laboris nisi ut aliquip ex ea commodo
consequat. Next is a theorem---numbering is specific to each type of entry, so theorems have a separate numbering from assumptions:
\begin{theorem}
Duis aute irure dolor in reprehenderit in voluptate velit esse
cillum dolore eu fugiat nulla pariatur. Excepteur sint occaecat cupidatat non
proident, sunt in culpa qui officia deserunt mollit anim id est laborum:
\begin{equation}
X^* =\iint_{0}^{\infty} \alpha(i) \cdot \mathcal{A}^2(i) + \mathbb{P}(X\mid Z(i))\,di\,dj,\end{equation}
integer vestibulum sapien nec velit varius, nec scelerisque nunc eleifend. Cras sollicitudin, justo sit amet tempus vehicula, ex ex vulputate nulla, non vestibulum quam tortor nec nisl. 
\end{theorem}

Lorem ipsum dolor sit amet, consectetur adipiscing elit. Sed aliquam magna vel urna ultrices, sed lacinia nulla mattis. Fusce non nunc nec est mollis malesuada. Here is another theorem---continuing the numbering:
\begin{theorem}
Consectetur adipiscing elit, sed do eiusmod tempor incididunt ut labore et dolore magna aliqua: $y(\gamma) \geq 3\pi + \cos(\vartheta)$.
\label{t:theorem1}\end{theorem}

\begin{proof}
Based on assumption~\ref{a:appe}, lorem ipsum dolor sit amet, consectetur adipiscing elit. Sed aliquam magna vel urna ultrices, sed lacinia nulla mattis. Fusce non nunc nec est mollis malesuada. Nullam dignissim nulla sit amet libero facilisis, eget fringilla libero sagittis. Suspendisse potenti. Vivamus fermentum consectetur ante, at rhoncus nisi tristique vel. Vivamus in est quis justo fermentum lacinia ac eu leo. Maecenas nec tempor nisi---as in theorem \ref{t:theorem1}.
\end{proof}

\section{More math}

Here is math and equations in the appendix---see equations \eqref{e:appendix1} and \eqref{e:experiments}. Temporibus autem quibusdam $\xi$ et aut officiis debitis aut rerum necessitatibus saepe eveniet ut et voluptates repudiandae sint et molestiae non recusandae $1-\gamma$. Itaque earum rerum hic $\mathcal{S}(\zeta^0)$ tenetur a sapiente delectus $\mathcal{B}^\theta$, ut aut reiciendis voluptatibus maiores alias consequatur aut perferendis doloribus asperiores repellat $\mathbb{V}^i$:
\begin{equation}
\mathbb{V}^r = (1-\gamma) \times 0 +\gamma S(z^*) v^s+\gamma [1-S(z^*)] \mathcal{V}^i-c.
\label{e:appendix1}\end{equation}

Ut enim ad minima veniam, quis nostrum exercitationem ullam corporis suscipit laboriosam, nisi ut aliquid ex ea commodi consequatur? Quis autem vel eum iure reprehenderit qui in ea voluptate velit esse quam nihil molestiae consequatur, vel illum qui dolorem eum fugiat quo voluptas nulla pariatur? 
\begin{equation}
\mathbb{E}(N(z)) = \frac{1}{1-\gamma F(z)}.
\label{e:experiments}\end{equation}
